% ============================================================
%  Stage-1 DBT fine-tuning results (Overleaf-ready)
%  Requires: \usepackage{booktabs}   % add to your preamble
%  Held-out DBT TEST/VAL AUROC. Backbone: 2.5D per-slice Swin-B (+ MIL).
%  Combined training data: BCS-DBT + DBT-2026.
%  Both rows are DBT fine-tuned; they differ only in initialization.
%  ImageNet->DBT = ImageNet-pretrained init then DBT fine-tune (backbone source of m1).
%  Scratch->DBT  = random init then DBT fine-tune (backbone source of mscratch).
%  Delta_init = (ImageNet->DBT) - (Scratch->DBT)  (gain from ImageNet initialization)
% ============================================================

% ---- Minimal standalone wrapper (delete if pasting into an existing paper) ----
\documentclass{article}
\usepackage{booktabs}
\begin{document}

\begin{table}[t]
  \centering
  \caption{Stage-1 DBT fine-tuning AUROC on the held-out DBT split
  (2.5D per-slice Swin-B + MIL, trained on combined BCS-DBT + DBT-2026).
  The ImageNet-only model (m0 source) has no DBT stage, hence no DBT score (---).
  $\Delta_{\mathrm{init}}=$(ImageNet$\rightarrow$DBT)$-$(Scratch$\rightarrow$DBT). Higher is better.}
  \label{tab:finetune_dbt}
  \begin{tabular}{lcc}
    \toprule
    Model (init $\rightarrow$ DBT fine-tune) & VAL AUROC & TEST AUROC \\
    \midrule
    ImageNet $\rightarrow$ DBT (m1 source)      & $\mathbf{0.951}$ & $\mathbf{0.929}$ \\
    Scratch $\rightarrow$ DBT (mscratch source) & $0.894$ & $0.851$ \\
    \midrule
    $\Delta_{\mathrm{init}}$ (m1 $-$ mscratch)  & $+0.057$ & $+0.078$ \\
    \bottomrule
  \end{tabular}
\end{table}

\end{document}
